\section{Results}

\subsection{Study design and self-supervised pretraining}

We pretrained \method{} on 169,120 human DNA methylation profiles represented in a shared
vocabulary of 49,156 CpG sites. Pretraining was label-free: neither chronological age nor
tissue, sex, dataset, or disease labels contributed to the objective. Each profile generated
two independently sampled views, each containing 50\% of its measured CpGs. A shared encoder
compressed each view into a 256-dimensional CLS representation, from which a decoder predicted
measured CpGs withheld from that view. A symmetric cross-view contrastive term encouraged the
two CLS representations from the same profile to agree while retaining discrimination between
profiles (Fig.~\ref{fig:study-design}).

We evaluated the pretrained encoder on the AltuMAge-derived collection, which contains 10,988
profiles and 21,368 measured CpGs before age-range filtering and duplicate removal. We used the
original train, validation, and test partition (7,416/1,308/2,264 profiles) only for frozen
representation diagnostics. The supervised age benchmark applied the prespecified age filter
and methylation-fingerprint deduplication, generated five stratified training/validation folds,
and retained a common fixed test set of 2,149 unique sample identifiers for every fold and both model
families. This separation is important: the 2,264-profile split supports in-distribution probes,
whereas all reported fine-tuning and paired-comparison metrics use the filtered and deduplicated
2,149-sample test set.

The pretraining corpus size was confirmed directly from the canonical data header
(169,120 profiles $\times$ 49,156 CpG sites) and partitioned 80/10/10 by a fixed permutation,
leaving 16,912 profiles that were never seen during pretraining and are used for the
reconstruction analyses below. This count should not be conflated with the 154,063 post-QC
figure reported for MethylGPT, whose curation pipeline differs.

\begin{figure}[H]
 \centering
 \includegraphics[width=\linewidth]{fig1_study_design_wced.pdf}
 \caption{\textbf{Study design and self-supervised pretraining of \method{}.}
 \textbf{a}, Study design. The unlabelled pretraining corpus is distinguished from the labelled
 AltuMAge evaluation collection, and the two downstream analysis populations are shown
 separately: the 10,988-profile partition used for frozen-representation probes and the
 filtered, deduplicated 2,149-sample test set used for supervised comparisons. Phenotype
 annotations were not used as supervision during pretraining. \textbf{b}, View construction.
 Two independent 50\% subsets of the measured CpGs of one profile form the two views; for each
 view, the complement (hatched) provides the reconstruction target, so the objective is scored
 only at measured CpGs absent from that view. \textbf{c}, The WCED objective. A shared encoder
 maps each view to a CLS representation; a decoder predicts the withheld CpGs from that
 representation alone, and a symmetric InfoNCE term aligns the two projected CLS embeddings.
 \textbf{d}, Token encoding. Each CpG token sums a learned site-identity embedding and a
 continuous $\beta$-value encoding, with a CLS token prepended to every partial profile.}
 \label{fig:study-design}
\end{figure}

\subsection{WCED learns view-consistent global representations}

WCED is designed to compress incomplete observations of a methylation profile into a shared
global representation. Across five independent random draws of two half-profile views for each
of 2,000 profiles, matched views had a mean cosine similarity of $0.984\pm0.0004$, compared
with $0.499\pm0.0002$ for unmatched profiles, and the correct paired view was retrieved for
73.8\% of queries (range 73.1--74.5\% across draws; chance 0.05\%). Retrieval reached 96.1\%
within the ten nearest candidates. The representation is therefore largely invariant to which
half of the methylome is observed while remaining specific enough to identify the individual
profile.

This invariance is nevertheless partly supported by shared input rather than by the
representation alone. Because the two views are drawn independently, they share approximately
half of their CpGs in expectation. Repeating the evaluation with strictly disjoint views---a
random partition in which the two halves share no site---left mean matched similarity almost
unchanged ($0.971\pm0.0005$ versus $0.984$) but reduced top-1 retrieval to
25.3\% (range 24.6--26.0\%). The dissociation is informative: average similarity is insensitive
to the change, whereas identifying the correct partner among 1,999 competitors depends on
narrow margins against the most similar profiles, which the small loss of positive similarity
is sufficient to erase. Cross-subset identity information is clearly retained---disjoint-view
retrieval remains roughly 500-fold above chance, and the correct partner falls within the ten
nearest candidates for 75.1\% of queries, comparable to top-1 performance under overlapping
views---but it is weaker than the overlapping-view figure implies. We note that the encoder was
trained with overlapping views, so disjoint evaluation is out of distribution and these values
are a conservative bound rather than an estimate of what the architecture achieves when trained
for that regime (Supplementary Fig.~S1).

A related question is how far the representation of a partial profile departs from that of the
complete profile, which bears on whether a sample assayed at fewer sites can be embedded
consistently. The embedding of a random 50\% CpG subset had a mean cosine similarity of 0.991
to the embedding of its own complete profile (shuffled-pairing control, 0.496), and a ridge age
probe fitted on complete-profile embeddings and applied without adaptation to subset embeddings
of held-out profiles retained $R^2=0.835$ against $R^2=0.863$ on complete profiles (paired
difference 0.028, 95\% CI 0.014--0.044; Supplementary Fig.~S2). The representation is therefore
substantially, though not perfectly, robust to which half of the measured sites is retained.
We note that these are random subsets rather than the probe set of a particular array
generation, so this does not by itself establish transfer across measurement platforms.

\subsection{WCED reconstructs CpGs withheld from the input}

We evaluated reconstruction on 5,000 profiles from the held-out portion of the pretraining
corpus, encoding each from a 50\% subset of its measured CpGs and scoring predictions only at
measured CpGs absent from that input ($1.12\times10^{8}$ evaluated positions). Predicted and
observed beta values were closely correlated at these withheld sites (Pearson $r=0.972$), and
the per-profile standardized reconstruction loss was 0.0525---closely reproducing the
checkpoint's own validation diagnostics (0.9713 and 0.0552) on profiles excluded from
pretraining (Fig.~\ref{fig:partial-view}c).

To test whether the decoder genuinely reads the sample representation rather than exploiting
population-level regularities, we compared three controls computed on the same profiles under
identical withheld masks. Reconstruction error was 0.0056 for the model, against 0.0306 for a
per-CpG mean predictor, 0.0505 for a decoder given standard normal vectors, and 0.0552 for a
decoder given CLS embeddings shuffled across the batch (Fig.~\ref{fig:partial-view}d). The
model therefore reduced error 5.4-fold relative to the per-CpG mean, even though the mean
predictor was computed on the evaluation cohort itself and is thus advantaged. Notably,
shuffled CLS embeddings produced slightly worse reconstruction than random vectors: supplying
a real but mismatched sample representation is more misleading than supplying no sample
information at all. Together with cross-view consistency, this shows that the CLS bottleneck
carries sample-specific information sufficient both to identify a profile and to predict
methylation values never supplied to that view.

\begin{figure}[H]
 \centering
 \includegraphics[width=\linewidth]{fig2_wced_consistency_reconstruction.pdf}
 \caption{\textbf{WCED learns consistent partial-view representations and reconstructs
 withheld CpGs.} \textbf{a}, Cosine similarity between CLS embeddings of two independently
 sampled 50\%-CpG views, for 40 representative profiles (heatmap; the bright diagonal marks
 matched views of the same profile) and summarised over all 2,000 profiles (violins; means
 annotated). \textbf{b}, Cross-view retrieval accuracy as a function of the number of
 candidates considered, against chance; the strictest point, top-1, is annotated. Panels a and
 b show one representative draw of two views per profile; values across five independent draws
 are given in the text and in Supplementary Fig.~S1. \textbf{c}, Observed
 versus predicted $\beta$ values at CpGs withheld from the encoder input, for 5,000 profiles
 from the held-out portion of the pretraining corpus (hexagonal binning, logarithmic colour
 scale; dashed line indicates identity). For legibility the panel displays a random subsample
 of $5.0\times10^{5}$ of the $1.12\times10^{8}$ evaluated positions; the reported correlation
 ($r=0.972$) was computed over all evaluated positions, not over the plotted subsample.
 \textbf{d}, Mean squared error at withheld CpGs for the model and three controls evaluated on
 the same profiles under identical withheld masks. Bars show the mean across profiles and
 vertical lines span the 10th--90th percentile across profiles; these are dispersion ranges,
 not confidence intervals.}
 \label{fig:partial-view}
\end{figure}

\subsection{Pretraining captures age and tissue variation}

We next asked whether self-supervised pretraining preserved information useful beyond its
reconstruction objective. A ridge regressor fitted to frozen pretrained CLS embeddings on the
original train-plus-validation profiles achieved $R^2=0.866$ and a median absolute error
(MedAE) of 5.17 years on the original 2,264-profile test split. Applying the same probe to
mean-pooled contextual CpG embeddings reduced performance to $R^2=0.795$ and MedAE 6.53 years.
The CLS space also had a higher effective rank than the mean-pooled space (18.12 versus 3.40)
and lower mean cosine anisotropy (0.500 versus 0.867). These findings indicate that the explicit
sample token retained broader, more age-informative variation than uniform pooling across CpG
tokens, despite receiving no age supervision during pretraining.

Biological information was not restricted to age. A linear classifier trained on frozen CLS
embeddings attained balanced accuracy 0.936 across 37 sufficiently represented tissue classes
(chance level 0.027). These probes establish that biologically associated structure is
accessible in the pretrained representation under an in-distribution sample split. They do not
show that the representation has separated biology from study or platform effects, because
profiles from the same study can occur on both sides of the probe split.

An archived nonlinear replica-head probe reached $R^2=0.922$ and MedAE 3.57 years, but we treat
it as exploratory: the current analysis script does not seed its NumPy minibatch shuffling and
does not reserve a separate split for head-hyperparameter selection. The deterministic ridge
probe is therefore the primary evidence for label-free age information.

\subsection{Age prediction is stable across folds}

We fine-tuned five models using different stratified training/validation folds and evaluated
every selected checkpoint on the same 2,149-sample test set (age range 0--103 years; mean
45.53 years). Across folds, mean MedAE was $3.169\pm0.040$ years, mean MAE was
$4.493\pm0.059$ years, and mean $R^2$ was $0.9316\pm0.0024$ (mean $\pm$ sample SD;
Table~\ref{tab:folds}). The narrow ranges across folds show that performance was not driven by
one favorable supervised partition. Because the test samples were shared, however, these
standard deviations measure training-partition stability and are not confidence intervals for
population performance.

Averaging the five predictions for each sample produced the primary \method{} estimate, with
MedAE 3.112 years, MAE 4.333 years, and $R^2=0.9365$. MethylGPT was aggregated in the same way
for the paired neural-model comparison below.

\begin{table}[H]
\centering
\caption{\textbf{Five-fold stability on the common held-out test set.}
Values were recomputed from saved sample-level predictions. Because every model was evaluated
on the same test samples, fold dispersion quantifies sensitivity to the training partition and
is not a population confidence interval.}
\label{tab:folds}
\begin{tabular}{lrrr}
\toprule
Model & MedAE (years) & MAE (years) & $R^2$ \\
\midrule
Fold 0 & 3.126 & 4.439 & 0.9321 \\
Fold 1 & 3.213 & 4.458 & 0.9337 \\
Fold 2 & 3.209 & 4.570 & 0.9284 \\
Fold 3 & 3.152 & 4.541 & 0.9298 \\
Fold 4 & 3.144 & 4.455 & 0.9338 \\
\midrule
Mean $\pm$ SD & $3.169\pm0.040$ & $4.493\pm0.059$ & $0.9316\pm0.0024$ \\
Fold ensemble & 3.112 & 4.333 & 0.9365 \\
\bottomrule
\end{tabular}
\end{table}

\subsection{\method{} improves age prediction over MethylGPT}

We performed a paired comparison using identical held-out samples, age labels, and fold
assignments for both model families. Fold-level predictions were averaged within sample, and
uncertainty was estimated by resampling the same test samples for both models. \method{}
reduced ensemble MedAE from 3.70 to 3.11 years (paired difference, 0.59 years; 95\% CI,
0.34--0.81) and MAE from 5.11 to 4.33 years, while increasing $R^2$ from 0.911 to 0.936.
Paired confidence intervals for all three endpoints excluded zero
(Fig.~\ref{fig:age-benchmark}).

We additionally evaluated a conventional linear reference using ElasticNetCV
\citep{Zou2005}. Hyperparameters were selected by internal cross-validation on the pooled
training and validation profiles, with the common test set withheld until final evaluation.
The resulting model used 1,036 non-zero coefficients among 21,368 CpGs and achieved MedAE
3.44 years, MAE 4.73 years, and $R^2=0.9322$. The \method{} ensemble therefore had better
point estimates than both ElasticNet and MethylGPT on this test set. Because ElasticNet was
evaluated as a single fitted model and its sample-level predictions were not retained, this
comparison is descriptive and is not assigned a paired confidence interval.

The benchmark is matched at the sample and split levels, but the pipelines differ in their
treatment of 1,760 CpGs outside the shared pretraining reference: \method{} retains their
measured values under a common unknown-CpG identity, whereas MethylGPT excludes them. The
result therefore establishes a predictive advantage for the implemented \method{} pipeline,
but does not attribute that advantage to WCED alone.

\begin{table}[H]
\centering
\caption{\textbf{Age-prediction performance on the common held-out test set.}
The first three rows are test-set point estimates. \method{} and MethylGPT are five-fold
prediction ensembles; ElasticNetCV is a single model selected by internal cross-validation.
The final row reports the paired MethylGPT-minus-\method{} difference and its 95\% bootstrap
confidence interval; positive error differences and a negative $R^2$ difference favor
\method{}.}
\label{tab:age-benchmark}
\small
\begin{tabular}{@{}lccc@{}}
\toprule
Model or comparison & MedAE (years) & MAE (years) & $R^2$ \\
\midrule
\method{} & 3.112 & 4.333 & 0.9365 \\
ElasticNetCV & 3.437 & 4.731 & 0.9322 \\
MethylGPT & 3.701 & 5.105 & 0.9110 \\
\midrule
MethylGPT $-$ \method{} & $0.589\ [0.338,\,0.815]$ &
$0.773\ [0.583,\,0.968]$ & $-0.0254\ [-0.0329,\,-0.0183]$ \\
\bottomrule
\end{tabular}
\end{table}

\begin{figure}[H]
 \centering
 \includegraphics[width=\linewidth]{fig3_age_benchmark_paired_comparison.pdf}
 \caption{\textbf{Stable age prediction and comparison with neural and linear references.}
 \textbf{a}, Test-set MedAE, MAE, and $R^2$ for the five independently fine-tuned \method{}
 folds; horizontal lines denote fold means. \textbf{b}, Ensemble-predicted versus chronological
 age for \method{} and MethylGPT on the same 2,149 held-out samples; each point is one
 sample and dashed lines indicate identity. \textbf{c}, Ensemble point estimates for
 \method{} and MethylGPT and the single-model ElasticNetCV reference. \textbf{d}, Sampling
 distributions of the paired difference between models over 10,000 sample resamples, with the
 central 95\% shaded; error reductions are MethylGPT-minus-\method{} differences, whereas the
 $R^2$ effect is \method{}-minus-MethylGPT, so positive values favor \method{} throughout. The
 vertical dashed lines mark no difference.}
 \label{fig:age-benchmark}
\end{figure}

\subsection{Genomic RoPE induces contextual CpG locality}

To test the representational effect of genomic RoPE, we trained a matched pilot pair on the
same fixed 5,000-profile subset. The two runs used the same architecture, objective, and
hyperparameters; the intentional difference was whether RoPE received genomic-rank positions
or sequential positions. For the genomic-RoPE model, mean cosine similarity of contextualized
representations was 0.9831 for adjacent CpGs and 0.8857 for random pairs, giving a locality gap
of 0.0974. Without genomic positions, the corresponding values were 0.9931 and 0.9621, a gap
of 0.0310. Thus, relative separation of adjacent from random CpGs was approximately 3.14-fold
larger with genomic positions. In the production model's static token-embedding table, adjacent
and random CpGs showed no positive locality gap, consistent with locality emerging through
contextual computation rather than simple clustering of identity embeddings.

This ablation supports a mechanistic interpretation but remains a pilot. It used one seed, and
the selected genomic checkpoint also had better reconstruction diagnostics than the non-genomic
checkpoint (normalized reconstruction loss 0.0906 versus 0.1585; Pearson correlation 0.9357
versus 0.9017). We therefore report the locality result in the Supplementary Results and do not
claim a downstream predictive advantage for genomic RoPE until replicated pretraining runs and
a matched downstream ablation are available.

\subsection{Pretrained representations retain tissue and study structure}

Although tissue was highly predictable from frozen CLS embeddings, study identity was even
more readily accessible: a linear probe attained balanced accuracy 0.956 across 85 eligible
datasets, compared with a chance level of 0.012. Two-dimensional projections likewise formed
many compact study-associated islands. This finding is not evidence of batch invariance.
Instead, it shows that the representation simultaneously retains biological tissue structure
and substantial cohort-specific structure. Tissue and dataset labels are also partially
confounded in public methylation collections, so the tissue-probe accuracy cannot be interpreted
as a study-independent biological metric.

For this reason, the main representation figure should pair the tissue result with the dataset
result rather than show tissue alone. UMAP panels are descriptive and belong in the
Supplementary Information; balanced-accuracy probes and cross-study tests provide the
quantitative evidence (Fig.~\ref{fig:representation-generalization}).

\subsection{Cross-study transfer of age representations remains limited}

We assessed whether age information in frozen CLS embeddings transferred to studies excluded
from probe fitting. Whereas a random-split ridge probe achieved $R^2=0.856$, the median
leave-one-dataset-out $R^2$ across 11 eligible studies was 0.261, with a median MedAE of
8.32 years; three studies yielded negative $R^2$. The broad variation between cohorts shows
that readily decodable age information within a mixed-study split does not ensure robust
transfer to an unseen study.

The failures are not explained by the obvious cohort descriptors. Neither the age dispersion
of the held-out cohort (Spearman $\rho=+0.28$, $P=0.40$) nor its sample size
($\rho=+0.26$, $P=0.43$) was associated with transfer success, and the failing studies overlap
the succeeding ones on both axes: one cohort of 118 profiles with an age SD of 10.5 years
transferred at $R^2=0.42$, whereas a larger cohort of 282 profiles with a wider age SD of
13.9 years transferred at $R^2=-4.63$ (Fig.~\ref{fig:representation-generalization}c). Cohort
size and age range therefore do not account for the heterogeneity, which instead points toward
study-specific biological or technical structure of the kind quantified above.

This is a diagnostic of the frozen representation, not a generalization test of the fully
fine-tuned age model. Differences in tissue, disease status, platform, and age distribution
also remain entangled with study identity. Claims of study-independent age prediction will
therefore require grouped study-held-out fine-tuning with cohort-level uncertainty; the
current analysis identifies cross-study transfer as an unresolved limitation.

\begin{figure}[H]
 \centering
 \includegraphics[width=\linewidth]{fig4_representation_cross_study_transfer.pdf}
 \caption{\textbf{Biological structure, study structure, and the limits of cross-study
 transfer.} \textbf{a}, Frozen linear age-probe performance using the pretrained CLS
 representation or mean-pooled contextual CpG embeddings. \textbf{b}, Balanced accuracy of
 frozen CLS linear probes for tissue and dataset/study identity; dark-grey dashed horizontal
 segments mark chance levels. \textbf{c}, Leave-one-dataset-out age-probe $R^2$ for the 11 studies meeting
 the prespecified size ($n\geq100$) and age-dispersion (SD $\geq10$ years) thresholds, plotted
 against the age spread of the held-out cohort, with marker area proportional to cohort size.
 Studies with negative $R^2$ are shown in red on a shaded background; the most extreme value is
 clipped for readability and annotated with its true value. Neither cohort descriptor predicts
 transfer success (Spearman $\rho$ inset). These analyses diagnose the frozen representation and
 do not constitute external validation of the fully fine-tuned model.}
 \label{fig:representation-generalization}
\end{figure}
